\documentclass[11pt]{article}
\usepackage[a4paper,margin=21mm]{geometry}
\usepackage[no-math]{fontspec}
\setmainfont{Latin Modern Roman}
\newfontfamily\CJKfont{Noto Serif CJK SC}
\XeTeXlinebreaklocale "zh"
\XeTeXlinebreakskip=0pt plus 1pt
\usepackage{amsmath,amssymb,amsthm,amscd,graphicx,color}
\usepackage{xurl}
\usepackage[unicode,hidelinks]{hyperref}
\allowdisplaybreaks
\setlength\emergencystretch{3em}
\numberwithin{equation}{section}


\numberwithin{equation}{section}

\newtheorem{Theorem}{Theorem}[section]
\newtheorem*{Theorem*}{Theorem}
\newtheorem{Corollary}[Theorem]{Corollary}
\newtheorem{Lemma}[Theorem]{Lemma}
\newtheorem{Proposition}[Theorem]{Proposition}
{ \theoremstyle{definition}
\newtheorem{Definition}[Theorem]{Definition}
\newtheorem{Note}[Theorem]{Note}
\newtheorem{Example}[Theorem]{Example}
\newtheorem{Remark}[Theorem]{Remark}
\newtheorem{Question}[Theorem]{Question}
}


\newtheorem{thmB}[Theorem]{Theorem B}
\newtheorem{thm}[Theorem]{Theorem}
\newtheorem{claim}[Theorem]{Claim}



\theoremstyle{plain}
\newtheorem*{namedthm}{\namedthmname}
\newcounter{namedthm}

\newtheorem{bigthm}{Theorem}
\renewcommand{\thebigthm}{\Alph{bigthm}}


\newcommand{\B}{\mathbb B}
 \newcommand{\D}{\mathbb D}
 \newcommand{\R}{\mathbb R}
 \newcommand{\Q}{\mathbb Q}
 \newcommand{\C}{\mathbb C}
 \newcommand{\PP}{\mathbb P}
 \newcommand{\N}{\mathbb N}
 \newcommand{\Z}{\mathbb Z}
 \newcommand{\Ox}{\mathcal O}

 \newcommand{\prv}{Proof: }
 \newcommand{\e}{\varepsilon}
 \newcommand{\vep}{\varepsilon}
 \newcommand{\capo}{{\rm Cap}_{\omega}}
 \newcommand{\capa}{{\rm Cap}}
 \newcommand{\capfp}{{\rm Cap}_{\varphi,\psi}}
 %\newcommand{\capp}{{\rm Cap}_{\phi}}
 \newcommand{\captheta}{{\rm Cap}_{\theta}}


 \newcommand{\f}{\varphi}
 \newcommand{\fm}{\varphi_{\min}}
 \newcommand{\p}{\psi}
 \newcommand{\s}{\sigma}
 \newcommand \g {\gamma}
 \newcommand \al {\alpha}
 \newcommand \de {\delta}
 \newcommand \la {\lambda}
 \newcommand \La {\Lambda}
 \newcommand \psh {{\rm PSH}}
 \newcommand \PSH {{\rm PSH}}
 \newcommand \SH {{\rm SH}}
 \newcommand \nhd {neighborhood }
 \newcommand \Om {\Omega}
 \newcommand{\id}{{\bf 1}}
 \newcommand \MA {{\rm MA}}
 \newcommand \Amp {{\rm Amp}}
 \newcommand \vol{{\rm Vol}}
 \newcommand \Ent{{\rm Ent}}
\newcommand{\Ric}{{\rm Ric}}
 \newcommand \cE{\mathcal E }
 \newcommand \Ec{\mathcal E }
 \newcommand \E{\mathcal E^1 (X,\omega) }
 \newcommand \cEc{\mathcal E_{\chi} (X,\omega) }
 \newcommand \mD{\mathcal D}
 \newcommand \mK{\mathcal K}
 \newcommand \mrm{\mathrm}
 \newcommand \sm{\setminus}
 \newcommand \es {\emptyset}
 \newcommand \Sub {\Subset}
 \newcommand \sub{\subset}
 \newcommand \nbd{neighborhood}
 \newcommand \lra{\longrightarrow}
 \newcommand \Lra{\Longrightarrow}
 \newcommand \lto{\longmapsto}
 \newcommand \mc {\mathcal}
 \newcommand \mb {\mathbb}
 \newcommand \we{\wedge}
 \newcommand \bwe{\bigwedge}
 \newcommand \ha {\Hat}
 \newcommand \ove {\overline}
 \newcommand \bd {\partial}

 \newcommand \cM {\mathcal{M}}

\begin{document}
\begin{center}\Large General-weight relative Moser–Trudinger estimate for positive-mass model potentials in a big class on a compact Kähler manifold\end{center}
\noindent\textbf{Stable ID:} 983\quad\textbf{Author:} Eleonora Di Nezza; Stefano Trapani; Antonio Trusiani\par
\noindent\textbf{Work:} Entropy for Monge--Amp\`ere Measures in the Prescribed Singularities Setting\par
\noindent\textbf{Source:} \url{https://sigma-journal.com/2024/039/}\par
\noindent\textbf{Version:} Official SIGMA 20 (2024) article 039, DOI 10.3842/SIGMA.2024.039; journal PDF and native TeX acquired 2026-09-30, exact bytes pinned by SHA256\par
\noindent\textbf{Rights:} CC BY 4.0. Authors retain copyright. \url{https://creativecommons.org/licenses/by/4.0/}. Reuse and typesetting adaptation; original language and mathematical content preserved.\par
\noindent\textbf{Difficulty (prior selection preserved):} {\CJKfont H3: The author constructs a nonlinear envelope using the integral of the nth root of an arbitrary admissible weight and its inverse, then compares non-pluripolar measures on the contact set. A mass contradiction determines the energy scale, and a separate convex-threshold argument yields uniform exponential integrability. These envelope and measure estimates are the substantive obligation beyond ordinary graduate integration inequalities.}\par
\medskip\noindent\textbf{Editorial boundary note (not author text):} Complete original Theorem3.5 and its full general-weight nonlinear-envelope, inverse-weight/contact-set, positive-mass contradiction, convex-threshold and uniform exponential-integrability proof. Every compact-Kahler/big-class, model-potential/positive-mass, continuous strictly increasing unbounded weight, relative full-mass/energy and sup-normalization definition is retained, including the full current truncation-energy monotonicity proof used in the initial reduction and the normalized volume/mass conventions. The precise cited standard envelope-support, non-pluripolar-measure comparison and uniform Skoda prerequisites remain formal author references; all new selected envelope and energy-scale arguments are fully human-authored and retained. Independent entropy/stability applications and their outcomes are excluded and not counted. Author literal differentiation formulas, coefficients and index wording are preserved without normalization or formula repair; no assistant proof or mathematical refereeing.\par
\medskip\hrule\medskip\noindent\textbf{Author text begins below.}\par
\setcounter{section}{1}
% Original source lines 224--444
\section{Preliminaries}\label{sec:prelim}

We recall results from (relative) pluripotential theory of big cohomology classes. We borrow notation and terminology from \cite{DDL6}.

Let $(X,\omega)$ be a compact K\"ahler manifold of dimension $n$ and let $\theta$ be a smooth closed $(1,1)$-form on $X$. A function $\varphi\colon X\rightarrow \mathbb{R}\cup\{-\infty\}$ is quasi-plurisubharmonic (qpsh) if it can be locally written as the sum of a plurisubharmonic function and a smooth function, and $\varphi$ is called $\theta$-plurisubharmonic ($\theta$\emph{-psh}) if it is qpsh and $\theta+{\rm dd}^c \varphi\geq 0$ in the sense of currents. Here, ${\rm d}$ and ${\rm d}^c$ are real differential operators defined as ${\rm d}:=\partial +\bar{\partial}$, ${\rm d}^c:=\frac{\rm i}{2\pi}\left(\bar{\partial}-\partial \right)$. We let $\psh(X,\theta)$ denote the set of $\theta$-psh functions that are not identically $-\infty$.
We also assume that $\{\theta\}$ is \emph{big}, i.e., there exists $\psi \in {\rm PSH}(X,\theta)$ such that $\theta +{\rm dd}^c \psi \geq \vep \omega$ for some small constant $\vep>0$. The current $T=\theta+{\rm dd}^c\psi$ is called K\"ahler current.

We say that a $\theta$-psh function $\varphi$ has \emph{analytic singularities} if there exists a constant $c>0$ such that locally on $X$,
\[
\varphi={c}\log\sum_{j=1}^{N}|f_j|^2+g,
\]
where $g$ is bounded and $f_1,\dots,f_N$ are local holomorphic functions.
We say that $\varphi$ has analytic singularities with \emph{smooth remainder} if moreover $g$ is smooth.

The \emph{ample locus} $\Amp(\theta)$ of $\{\theta\}$ is the complement of the \emph{non-K\"ahler locus}
\[E_{\rm nK}(\theta)= \bigcap_{T \ {\text{K\"ahler current}}} \{x\in X \colon \nu(T,x)>0 \},\]
where $\nu(T,x)$ is the Lelong number of $T$ at the point $x$.
The ample locus $\Amp(\theta)$ is a Zariski open subset, and it is nonempty \cite{Bou04}. Also we note that $\Amp(\theta)$ only depends only the cohomology class $\{ \theta\}$.

If $\varphi$ and $\varphi'$ are two $\theta$-psh functions on $X$, then $\varphi'$ is said to be \emph{less singular} than $\varphi$, i.e., $\varphi \preceq \varphi'$, if they satisfy $\varphi\le\varphi'+C$ for some $C\in \mathbb{R}$.
We say that $\varphi$ has the same singularity as~$\varphi'$, i.e., $\varphi \simeq \varphi'$, if $\varphi \preceq \varphi'$ and $\varphi' \preceq \varphi$. The latter condition is easily seen to yield an equivalence relation, whose equivalence classes $[\varphi]$, $\varphi \in \psh(X, \theta)$, are called \emph{singularity types}.

A $\theta$-psh function $\varphi$ has \emph{minimal singularity type} if it is less singular than any other $\theta$-psh function. Such $\theta$-psh functions with minimal singularity type always exist, one can consider for example $V_\theta:=\sup \{ \varphi \ \theta\text{-psh},\, \varphi\le 0\text{ on } X \}$. Trivially, a $\theta$-psh function with minimal singularity type is locally bounded in $\Amp(\theta)$. It follows from \cite[Theorem~1.1]{DNT23} that $V_{\theta}$ is $C^{1, \bar{1}}$ in the ample locus ${\rm Amp}(\theta)$.

Given
\[\theta^1+{\rm dd}^c\varphi_1, \ \dots, \ \theta^p+{\rm dd}^c \varphi_p\] positive $(1,1)$-currents, where $\theta^j$ are closed smooth real $(1,1)$-forms, following the construction of Bedford--Taylor \cite{BT87} in the local setting, it has been shown in~\cite{BEGZ10} that the sequence of currents
\[
{\bf 1}_{\bigcap_j\{\varphi_j>V_{\theta_j}-k\}}\big(\theta^1+{\rm dd}^c \max(\varphi_1, V_{\theta_1}-k)\big)\wedge\dots\wedge (\theta^p+{\rm dd}^c\max(\varphi_p, V_{\theta_p}-k))
\]
is non-decreasing in $k$ and converges weakly to the so called \emph{non-pluripolar product}
\[%\label{eq: BEGZ10_def}
\big\langle \theta^1_{\varphi_1 } \wedge\dots\wedge\theta^p_{\varphi_p}\big\rangle .
\]
In the following, with a slight abuse of notation, we will denote the non-pluripolar product simply by $\theta^1_{\varphi_1 } \wedge\dots\wedge\theta^p_{\varphi_p}$.
When $p =n$, the resulting positive $(n,n)$-current is a Borel measure that does not charge pluripolar sets. Pluripolar sets are Borel measurable sets that are locally contained in the $-\infty$ locus of psh functions. As a consequence of \cite[Corollary 2.11]{BBGZ13} for any pluripolar set $A$, there exists $\psi\in \psh(X,\theta)$ such that $A\subset \{\psi=-\infty\}$.

For a $\theta$-psh function $\varphi$, the non-pluripolar product $\theta_\varphi^n$ is said to be the \emph{non-pluripolar Monge--Amp{\`e}re measure} of $\varphi$.

The volume of a big class $\{ \theta\}$ is defined by
\[
{\rm Vol}(\theta):= \int_{{\rm Amp}(\theta)} \theta_{V_\theta}^n.
\]
For notational convenience, we simply write ${\rm Vol}(\theta)$, but keeping in mind that the volume is a~cohomological constant. Note that ${\rm Vol}(\theta)>0$ as $\{\theta\}$ is big (cf.\ \cite{BEGZ10}).

A $\theta$-psh function $\varphi$ is said to have \emph{full Monge--Amp\`ere mass} if
\[
\int_X \theta_\varphi^n={\rm Vol}(\theta),
\]
and we then write $\varphi\in \mathcal{E}(X,\theta)$. By \cite[Theorem 1.16]{BEGZ10}, the set $\mathcal{E}(X,\theta)$ strictly contains the set of $\theta$-psh functions with minimal singularity type.

An important property of the non-pluripolar product is that it is local with respect to the plurifine topology (see \cite[Corollary~4.3]{BT87}, \cite[Section~1.2]{BEGZ10}).
This topology is the coarsest such that all qpsh functions on $X$ are continuous. For convenience, we record the following version of this result for later use.

\begin{Lemma} \label{lem: plurifine}
Fix closed smooth real big $(1,1)$-forms $\theta^1,\dots,\theta^n$. Assume that $\varphi_j$, $\psi_j$, $j=1,\dots,n$ are $\theta^j$-psh functions such that $\varphi_j =\psi_j$ on $U$, an open set in the plurifine topology. Then
\[
\id_{U} \theta^1_{\varphi_1} \wedge \dots \wedge \theta^n_{\varphi_n} = \id_{U} \theta^1_{\psi_1} \wedge\dots \wedge \theta^n_{\psi_n}.
\]
\end{Lemma}
Lemma \ref{lem: plurifine} will be referred to as the \emph{plurifine locality property}. We will often work with sets of the form $\{u<v\}$, where $u$, $v$ are quasi-psh functions. These are always open in the plurifine topology.

\subsection{Envelopes and model potentials}
Giving $f$ an upper semi-continuous function (usc for short) on $X$, we define the envelope of $f$ in the class ${\rm PSH}(X,\theta)$ by
\[
P_{\theta}(f) := \sup \{u\in {\rm PSH}(X,\theta) \colon u\leq f\},
\]
with the convention that $\sup\varnothing =-\infty$. Observe that $P_{\theta}(f)\in {\rm PSH}(X,\theta)$ if and only if there exists some $u\in {\rm PSH}(X,\theta)$ lying below $f$. Note also that $V_{\theta}=P_{\theta}(0)$, and that $P_{\theta}(f+C)= P_{\theta}(f)+C$ for any constant $C$.

 In the particular case $f=\min(\psi,\phi)$ for $\psi$, $\phi$ usc functions, we denote the envelope as
$ P_\theta(\psi,\phi):=P_\theta(\min(\psi,\phi))$.
We observe that $ P_\theta(\psi,\phi)= P_\theta(P_\theta(\psi),P_\theta(\phi))$, so without loss of generality, we can assume $\psi$, $\phi$ are two $\theta$-psh functions.

Starting from the ``rooftop envelope'' $ P_\theta(\psi,\phi)$, we introduce
\[P_\theta[\psi](\phi) := \big(\lim_{C \to \infty}P_\theta(\psi+C,\phi)\big)^*.\]
It is easy to see that $P_\theta[\psi](\phi)$ only depends on the singularity type of $\psi$. When $\phi = V_\theta$, we will simply write $P_\theta[\psi]:=P_\theta[\psi](V_\theta)$, and we refer to this potential as the \emph{envelope of the singularity type} $[\psi]$.

Since $\psi - \sup_X \psi \leq P_\theta[\psi]$, we have that $[\psi] \leq [P_\theta[\psi]]$ and typically equality does not happen. When $[\psi] = [P_\theta[\psi]]$, we say that $\psi$ has \emph{model singularity type}. In the (more particular) case~${\psi = P_\theta[\psi]}$, we say that $\psi$ is a \emph{model potential}.

It is worth to mention that given any $\theta$-psh function $\psi$ with positive mass, the associated envelope $P_\theta[\psi]$ is in fact a model potential \cite[Theorem~3.14]{DDL6}.
Also, we recall that by \cite[Remark~3.4]{DDL6}, we know that $\int_X \theta_\psi^n = \int_X \theta_{P_\theta[\psi]}^n$.

From now on (unless otherwise stated), $\phi$ will denote a model potential with positive mass, i.e., $\int_X \theta_\phi^n>0$. We say that a $\theta$-psh function $\varphi$ has $\phi$-relative minimal singularities if $\varphi \simeq \phi$.

\begin{Definition}
Given a model potential $\phi$, the relative full mass class $\mathcal{E}(X,\theta,\phi)$ is the set of all $\theta$-psh functions $u$ such that $u$ is more singular than $\phi$ and $\int_X \theta_u^n=\int_X \theta_{\phi}^n$.
\end{Definition}

\begin{Proposition}\label{mass}
Assume $\{ \theta \}$ is a big cohomology class, and $0 \leq m \leq V$ with $V = \int_X \theta_{V_\theta}^n$. Then there exists a model potential $ u \in \PSH(X,\theta)$ such that $\int_X \theta_u^n = m$.

Moreover, $V_\theta$ is the only model potential with Monge--Amp\`ere mass~$V$, whereas there are infinitely many model potentials with Monge--Amp\`ere mass $m < V$.
\end{Proposition}

\begin{proof} For the first statement, we need to treat the $1$-dimensional case separately. When $n=1$, the big cohomology class $\{\theta\}$ is actually K\"ahler, hence we can work with a K\"ahler form $\omega\in \{\theta\}$ with $\int_X \omega = V$ as a reference form.
Fix $a\in X$, then the measures $\omega$ and $V \delta_a$ define cohomologus closed positive $(1,1)$-currents. Hence, there exists $\psi_a \in \SH(X,\omega)$ such that $\omega + {\rm dd}^c \psi_a = V \delta_a$ and~$\sup_X \psi_a = 0$.
Note that on a fixed local holomorphic coordinate chart centered in $a$ we have that~$\psi_a $ writes as the sum of $ V \log(|z|)$ and a smooth function. Thus, $ \psi_a(a) = - \infty$. Observe moreover that the non-pluripolar product $( \omega + {\rm dd}^c \psi_a)$ is given by the product of the current~${\omega + {\rm dd}^c \psi_a}$ with the characteristic function $ {\bf 1}_{ \{ \psi_a > - \infty \} }$. This implies that $\omega + {\rm dd}^c \psi_a$ has mass zero. Since for $t \in [0,1]$, we have $ \omega +t {\rm dd}^c \psi_a \geq (1-t) \omega \geq 0$, the function
\[m(t):= \int_X (\omega+t{\rm dd}^c\psi_a)\in [0, V]\]
 defined on $[0,1]$ is continuous and $m(0)=V$, $m(1)=0$, we arrive at the conclusion.

Assume now $n\geq 2$. Let us consider $\pi \colon Y \to X$ the blow up of $X$ at a point $p \in X$ with exceptional divisor $E$. Let $[E]$ be the current of integration along $E$.

Note that the cohomology class $\{ \pi^*\theta \} $ is big and so is $\{\pi^*(\theta) - t [E]\}$ for $ |t| < \varepsilon$ with~$\varepsilon$ small enough.
Observe that the pseudoeffective cone of a compact K\"ahler manifold does not contain lines. In fact, let $\alpha $ and $\gamma$ be smooth closed real $(1,1)$-forms and assume that
${\{ \alpha \} + t \{ \gamma \}}$ is a pseudoeffective class for all
$t \in \mathbb{R}$.
Then for $k \in { \N}$, \smash{$ \frac{ \{ \alpha \} }{k}+ \{ \gamma \} = \frac{1}{k}( \{ \alpha \}+ k\{ \gamma \})$} and \smash{$\frac{ \{\alpha \}}{k}+ \{-\gamma \} = \frac{1}{k}( \{ \alpha \} + k\{-\gamma \})$} are pseudoeffective, but the pseudoeffective cone is closed, then letting $k\rightarrow +\infty$, we get that both
$\{ \gamma \}$ and $\{ - \gamma \} $ are pseudoeffective classes, meaning that~${\{ \gamma \} = 0}$.

Therefore, for large $t$, $\{\pi^*(\theta) - t [E] \}$ is no longer big. Let
\[a = \sup \{ t \in \mathbb{R}, \ t > 0 \colon \mbox{ the class $\{ \pi^*(\theta) - t[E] \}$ is big} \}.
\]
Then $ 0< a < +\infty $ and
$\{\pi^*\theta - a [E]\}$ is pseudoeffective but not big.
In other words, if $\gamma$ is a~smooth form representing
$\{ \pi^*\theta - a [E] \}$, then
\[\int_Y (\gamma+{\rm dd}^c V_\gamma)^n = 0. \]

Now, $\gamma+{\rm dd}^c V_\gamma + a[E]$ is a positive $(1,1)$-current and, since the non-pluripolar product does not put mass of analytic subsets, we also get
\[\int_Y ( \gamma +{\rm dd}^c V_\gamma + a[E])^n = 0. \]

On the other hand, $\{\gamma+{\rm dd}^c V_\gamma + a[E]\}= \{ \pi^*\theta \}$, hence by \cite[Proposition 1.2.4]{BouThesis}, there exists~${u_1\in \PSH(X,\theta)}$ such that
$\pi^*\theta_{u_1} = \gamma +{\rm dd}^c V_\gamma + a[E]$.
In particular, $ 0=\int_Y (\pi^*\theta_{u_1})^n =\int_X \theta_{u_1}^n $.
Therefore, since the function
\[m(t) := \int_X ( \theta + (t u_1 + (1-t) V_\theta))^n \in [0, V] \]
defined in $[0,1]$ is continuous, $m(0)=V$, $m(1)=0$, then for $0 \leq m \leq V$ there exists $t_0\in [0,1]$ such that $\int_X (\theta + {\rm dd}^c ( t_0 u_1 + (1-t_0) V_\theta)^n = m$. On the other hand, by \cite[Remark 3.4]{DDL6}, $u_t := tu_1 + (1-t)V_\theta$ and $P[u_t]$ have the same mass for every $t \in [0,1]$.

For the last statement, we observe that, since $P[u_t] \leq V_\theta$ for all $t \in [0,1]$, by \cite[Theorem~3.14]{DDL6}, $V_\theta$~is the only model potential with Monge--Amp\'ere mass $V$. Now, by~\cite[Corollary~1.18]{BouThesis}, the Lelong number of the function~$u_1$ at~$p$ is strictly positive. Therefore, for $0 < t \leq 1$ also the Lelong number of $u_t$ is strictly positive at~$p$. By \cite[Lemma~5.1]{DDL6}, so is the Lelong number of~$P[u_t]$. Therefore, fixing $0 \leq m < V$ and varying the point $p$, we obtain infinitely many distinct model potentials all of the same Monge--Amp\`ere mass $m$.
\end{proof}

As pointed out in \cite{GZ07}, and then in the relative setting in \cite{Gup23}, it is natural to consider weighted subspaces of $\mathcal E(X,\theta,\phi)$.

A \emph{weight} is a continuous strictly increasing function $\chi\colon [0,+\infty) \rightarrow [0,+\infty)$ such that $\chi(0)=0$ and $\chi(+\infty)=+\infty$. Denote by $\chi^{-1}$ its inverse function, i.e., such that $\chi\big(\chi^{-1}(t)\big)= t$ for all~${t\geq 0}$.

We fix $\phi$ a model potential and we let $\mathcal{E}_{\chi}(X,\theta,\phi)$ denote the set of all $u\in \mathcal{E}(X,\theta,\phi)$ such that
\[
E_{\chi}(u,\phi):=\int_X \chi(|u-\phi|) \theta_u^n <\infty.
\]
When $\phi=V_{\theta}$, we denote $\mathcal{E}(X,\theta)=\mathcal{E}(X,\theta,V_{\theta})$, $\mathcal{E}_{\chi}(X,\theta)=\mathcal{E}_{\chi}(X,\theta,V_{\theta}$) and $E_{\chi}(u)=E_{\chi}(u,V_{\theta})$.

Compared to \cite{GZ07}, we have changed the sign of the weight, but the weighted classes are the same.

Also, in the special case $\chi(t)=t^p$, $p>0$, we simply denote the relative energy class with~${\mathcal{E}^p(X,\theta,\phi)}$
and the corresponding relative energy $E_p(u,\phi)$.

These weighted Monge--Amp\`ere energy classes covers the all class $\mathcal{E}$, i.e.,

\begin{Lemma}\label{energy union} We have that
$\mathcal{E}(X,\theta,\phi)=\bigcup_{\chi}\mathcal{E}_\chi(X,\theta,\phi)$.
\end{Lemma}

\begin{proof}We give here the proof for reader's convenience.
The inclusion $ \supseteq $ clearly follows from the definition. In order to prove the other inclusion we need to show that, given $\varphi \in \mathcal{E}(X,\theta,\phi)$ normalized with $\sup_X \varphi =-1$, there exists a weight function $\chi$ such that $\int_X \chi (\phi-\varphi) {\rm d}\mu < +\infty$, where $\mu:=\theta_\varphi^n$. Also, for convenience we normalize~$\mu$ to have $\mu(X)=1$.

Set \smash{$\psi(t):= \frac{1}{\mu(\{\varphi <\phi -t\})}$}, and observe that $\psi\colon\R^+\rightarrow \R^+$ is an increasing function such that $\psi(0)=1$ and $\lim_{t\rightarrow +\infty} \psi(t) =+\infty$. The latter property follows from the definition of the non-pluripolar product. We also note that, since $t\rightarrow \mu(\{\varphi <\phi -t\})$ is a decreasing function, its discontinuity locus is at most countable. In particular, $\psi$ is continuous for all $t\in \R^+\setminus J$.

We now consider $h\colon \R^+\rightarrow \R^+$ such that $h> 0$ satisfying:
\[{\rm (a)} \ \int_1^{+\infty} h(s) {\rm d}s <+\infty, \qquad {\rm (b)} \ \int_1^{+\infty} {h(s)} \psi(s) {\rm d}s =+\infty.\]
 Such a function $h$ exists as we show below. We choose a sequence of distinct points $\{t_k\}_{k\in \N}\in \R^+$ such that $\psi(t_k) \geq k$. Such a sequence exists since $\psi$ is increasing and $\psi\rightarrow +\infty $ as $t$ goes to~$+\infty$. Also, we can take $t_1=1$ since $\psi(1)\geq \psi(0)=1$. We then set $I_k:=[t_{k}, t_{k+1}[$ and we define $h$ to be the piece-wise continuous function constant on each $I_k\colon h|_{I_k}= k^{-2} (t_{k+1}-t_k)^{-1} $. Then
 \[\int_1^{+\infty} h(s) {\rm d}s= \sum_{k\geq 1} \frac{1}{k^2} <+\infty \qquad \text{and}\qquad \int_1^{+\infty}h(s)\psi(s) {\rm d}s\geq \sum_{k\geq 1}\frac{1}{k} =+\infty.\]

We then define the function $\chi\colon \R^+\rightarrow \R^+$ as \smash{$\chi(t):= \int_0^ t\frac{h(s)}{\mu(\{\varphi <\phi -s\})} {\rm d}s$}. We infer that the fundamental theorem of calculus holds almost everywhere. Indeed, for $t\notin J\cup \{t_k\}_k$, given $\varepsilon>0$ (small enough), the mean value theorem ensures that there exists $c_\varepsilon \in (t, t+\varepsilon)$ such that
\[\frac{\chi (t+\varepsilon)-\chi (t)}{\varepsilon} = \frac{1}{\varepsilon} \int_t^{t+\varepsilon} h(s) \psi(s) {\rm d}s = h(c_\varepsilon) \psi(c_\varepsilon).\]
Sending $\varepsilon\to 0$, and using the continuity of $h\cdot \psi$ on $\R^+\setminus (J\cup \{t_k\}_k)$, we get that $(\chi'(t))^+=h(t)\psi(t)$.

The same arguments work for left derivative, giving $(\chi'(t))^+=(\chi'(t))^-=\chi'(t)=h(t)\psi(t)$ almost everywhere, as we wanted.
By definition, $\chi$ is strictly increasing (since $h> 0$) and $\chi(+\infty)=+\infty$, thanks to condition (b). Moreover, setting $\varepsilon:=\chi (1)>0$, we have
\begin{align*}
\int_X \chi (\phi-\varphi) {\rm d}\mu = & \int_{\varepsilon}^\infty \mu(\{ \chi (\phi-\varphi) >t\}) {\rm d}t=\int_{\varepsilon}^\infty \mu(\{ \phi-\varphi >\chi^{-1}(t)\}) {\rm d}t\\
=& \int_{1}^\infty \chi'(s) \mu(\{ \phi-\varphi >s\}) {\rm d}s= \int_{1}^\infty h(s) {\rm d}s <+\infty,
\end{align*}
thanks to condition (a). This concludes the proof.
\end{proof}

We conclude here with a simple lemma needed in the following that we could not find in this precise form in the literature:

\begin{Lemma}\label{energy increases}
 Let $u\in \mathcal{E}_{\chi}(X,\theta,\phi)$ and $u_j:=\max (u,\phi -j)$. Then $u_j \in \mathcal{E}_{\chi}(X,\theta,\phi)$ and $E_{\chi}(u_j,\phi) \nearrow E_{\chi}(u,\phi)$.
\end{Lemma}
\begin{proof}
 For $j\in \N$, by Lemma \ref{lem: plurifine}, we have $\id_{\{u>\phi -j\}} \theta_{u_j}^n= \id_{\{u>\phi-j\}} \theta_{u}^n$. Since $\int_X \theta_{u}^n=\int_X \theta_{u_j}^n=\int_X \theta_\phi^n$, we obtain
\begin{align*}
0\leq \limsup_{j\to +\infty} \chi(j)\int_{\{u\leq \phi-j\}}\!\theta_{u_j}^n = \limsup_{j\to +\infty} \chi(j)\int_{\{u\leq \phi-j\}}\!\theta_{u}^n
\leq \limsup_{j\to +\infty} \int_{\{u\leq\phi-j\}}\! \chi (\phi-u)\theta_{u}^n =0.
\end{align*}
Now,
\begin{align*}
\int_X \chi (|\phi-u_j|) \theta_{u_j}^n &= \int_{\{ u > \phi-j\}} \chi (|\phi-u_j|) \theta_{u_j}^n+ \int_{\{ u \leq \phi-j\}} \chi (\phi-u_j) \theta_{u_j}^n \\
 &= \int_{\{ u > \phi-j\}} \chi (|\phi-u|) \theta_{u}^n+ \int_{\{ u \leq \phi-j\}} \chi (\phi-u_j) \theta_{u_j}^n.
\end{align*}
The conclusion follows letting $j\rightarrow +\infty$.
\end{proof}

\section{Entropy}
 We recall that given two positive probability measures $\mu$, $\nu$, the relative entropy $\Ent(\mu, \nu)$ is defined as
\[
\Ent(\mu, \nu) := \int_X \log\left (\frac{{\rm d}\mu}{{\rm d}\nu}\right) {\rm d}\mu,
\]
if $\mu$ is absolutely continuous with respect to $\nu$, and $+\infty$ otherwise.

\begin{Remark}
Let $\mu$, $\nu$ positive probability measure with $\mu := f \nu $ absolutely continuous with respect to $\nu$. Then $\Ent(\mu,\nu) < + \infty$ if and only if $f\log f \in L^1(X,\nu)$, in fact if $f < 1$, $f \log f$ is bounded, and if $f \geq 1$, $f \log f \geq 0$. \label{L^1}
\end{Remark}

Once and for all, we normalize the K\"ahler form $\omega$ such that $\int_X \omega^n=1$. We consider $u\in \psh(X, \theta)$ such that $\theta_u^n=f \omega^n$, $0 \leq f$ and $m_u:=\int_X \theta_u^n >0$. Then $ u\in \mathcal{E}(X, \theta, \phi)$ for the model potential $\phi=P_\theta[u]$ \cite[Theorem 1.3]{DDL2}, and $m_u^{-1}\theta_u^n$ is a probability measure. We then define the $\theta$\emph{-entropy} of $u$ as
\begin{align}
 {\rm Ent}_\theta(u)&:= {\rm Ent}\big( m_u^{-1}\theta_u^n, {\omega^n}\big)=\int_X \log \left( \frac{m_u^{-1}\theta_u^n}{\omega^n} \right) m_u^{-1}\theta_u^n\nonumber\\
 &= m_u^{-1} \int_X f\log f \omega^n-\log m_u.\label{eq: entropy}
 \end{align}
By Jensen’s inequality, we have ${\rm Ent}_\theta(u)\geq 0$. Also, observe that the definition of the $\theta$-entropy does depend on the chosen volume form $\omega^n$ but its finiteness does not.

Also, the expression in~\eqref{eq: entropy} coincides with the definition of entropy in \cite{DGL20} when $P[u]=V_\theta$, i.e., when $u\in \mathcal{E}(X, \theta)$. The definition in~\eqref{eq: entropy} is indeed a generalisation that allows to consider any $\theta$-psh function not necessarily of full mass.

More generally, given two $\theta$-psh functions $u$, $v$ with $m_u, m_v>0$ we define
\[{\rm Ent}_\theta(u,v) :={\rm Ent}\big( m_u^{-1}\theta_u^n, m_v^{-1}\theta_v^n\big). \]
Also, if no confusion can arise, we simply write ${\rm Ent} (u) $ and ${\rm Ent} (u,v) $.

\begin{Definition}\label{def: entropy}
 We say that $u\in \psh(X, \theta)$ with $m_u>0$ has \emph{finite $\theta$-entropy} if ${\rm Ent}_\theta(u)<+\infty$. We denote by ${\rm Ent}(X, \theta)$ the set of all $\theta$-psh functions having finite $\theta$-entropy.
\end{Definition}
 By \eqref{eq: entropy}, ${\rm Ent}_\theta(u)<+\infty$ if and only if $\theta_u^n=f\omega^n$ and $\int_X f\log f\omega^n<+\infty$ or equivalently $\int_X (f+1)\log (f+1)\omega^n<+\infty$.

\setcounter{Theorem}{4}
% Original source lines 506--619
\subsection{Entropy and energy}
It was proved in \cite{DGL20, DNL21} that \[\Ent(X,\theta)\cap \mathcal{E}(X,\theta) \subset \cE^{\frac{n}{n-1}}(X,\theta).\] As will be shown below, one can extend these results to the case of prescribed singularities.

We start with an integrability result of Moser--Trudinger type for general weights:

\begin{Theorem}
 	\label{thm: rel Moser--Trudinger inequality weight}
 	Let $\phi$ be a model potential with $m_\phi>0$. Let $\chi_1\colon \R^+\rightarrow \R^+ $ be a weight. Let \smash{$\chi_2(t) := \int_0^t \chi_1(s)^{\frac{1}{n}} {\rm d}s$}. Then there exist $c>0$, $C>0$ depending on $X$, $\theta$, $n$, $m_\phi$ and $\omega$ such that, for all $\varphi\in \mathcal{E}_{\chi_1}(X,\theta, \phi)$ with $\sup_X \varphi=-1$, we have
 	\begin{equation}
 	\int_X \exp \big( c E_{\chi_1}(\varphi, \phi)^{-\frac{1}{n}} \chi_2(\phi-\varphi)\big)\omega^n \leq C.
 	\label{eq: Moser} \end{equation}
 \end{Theorem}

 \begin{proof}
 We first note that if we replace $\chi_1$ by $\alpha \chi_1$ with $\alpha$ positive constant the left-hand side of~\eqref{eq: Moser} does not change, so we may assume $\chi_1(1) = 1$. 	
We claim that it suffices to prove the above inequality for $\varphi$ with relative minimal singularities, i.e., $[\varphi]=[\phi]$. Indeed, given~${\varphi\in \psh(X, \theta)}$, $\varphi_j:=\max(\varphi, \phi-j)$ has the same singularity type as $\phi$. This would mean that
 \[\int_X \exp \big( c E_{\chi_1}(\varphi_j, \phi)^{-\frac{1}{n}} \chi_2(\phi-\varphi_j)\big)\omega^n \leq C.\]
Moreover, it follows from Lemma \ref{energy increases} that $ E_{\chi_1}(\varphi_j, \phi) \nearrow E_{\chi_1}(\varphi, \phi) $ as $j\rightarrow +\infty$. Fatou's lemma will then give the desired estimate.

Thus assume that $\varphi$ has relative minimal singularities.
Let $\psi=-a\chi_2(\phi-\varphi)+\phi$ where $a>0$ is a small constant to be suitably chosen. Then define $u=P_\theta(\psi)$ and $v=-\gamma_2 (\phi-u) +\phi$, where~${\gamma_2 \colon \mathbb{R}^+ \rightarrow \mathbb{R}^+}$ denotes the inverse function of $a\chi_2$, which is concave and increasing. We observe that $u$ is a $\theta$-psh function with the same singularity type as $\phi$ and that
\[v=-\gamma_2 (\phi-u) +\phi \leq -\gamma_2(\phi-\psi) +\phi= -\gamma_2(a\chi_2(\phi-\varphi)) +\phi = \varphi\] with equality on the contact set $\mathcal{C}=\{u=\psi\}$.

A simple computation gives
 \begin{flalign*}
 \theta +{\rm dd}^c v & = \gamma_2'(\phi-u) \theta_u + (1-\gamma_2'(\phi-u)) \theta_{\phi} - \gamma_2''(\phi-u) d(\phi-u) \wedge d^c (\phi-u)\\
&\geq \gamma_2'(\phi-u) \theta_u+ (1-\gamma_2'(\phi-u)) \theta_{\phi},
 \end{flalign*}
where in the above we used that $\gamma_2$ is concave. We consider the set $G:=\{\gamma_2'(\phi-u)< 1\}$. We are going to show by contradiction that for a suitable choice of $a$ we have $G \neq X$. So, we assume $G= X$.
It then follow that $v$ is $\theta$-psh and, by construction, we infer that $v$ has the same singularity type as $\phi$.

Recall that $\sup_X \varphi=\sup_X(\varphi-\phi)$ as $\phi$ is a model potential \cite[Lemma 3.5]{DDL6}, hence $\phi-\varphi \geq 1$. This implies that $\chi_1$ is increasing and $\chi_1(1) = 1$ giving that
 \[E_{\chi_1}(\varphi, \phi) = \int_X \chi_1(\phi-\varphi) \theta_{\varphi}^n \geq \chi_1(1)\int_X \theta_{\varphi}^n= \int_X \theta_{\phi}^n=m_\phi.\]

By \cite[Theorem 2.7]{DDL6}, the non-pluripolar Monge--Amp\`ere measure $(\theta+{\rm dd}^c u)^n$ is supported on $\mathcal{C}$, hence
\begin{align*}%\label{eq: DGL contact weight}
	(\theta+{\rm dd}^c v) &\geq \gamma_2'(\phi-u) (\theta+{\rm dd}^c u)
		 = (a\chi_2'(\phi-\varphi))^{-1} (\theta+{\rm dd}^c u)\quad \text{on $\mathcal{C}$}.
\end{align*}
The last identity follows from the fact that, since $(a\chi_2)'(\gamma_2(t))\gamma_2'(t)=1$ and $\gamma_2(\phi-u)=\phi-\varphi$ in $\mathcal{C}$, we have $\gamma_2'(\phi-u)= (a\chi_2'(\phi-\varphi))^{-1}$ in $\mathcal{C}$.

It follows from \cite[Lemma 5.19]{DDL6} that the above inequality between positive currents implies an inequality between the non-pluripolar measures (observe that this is not trivial since $(a\chi_2'(\phi-\varphi))^{-1} (\theta+{\rm dd}^c u)$ is not closed).
Thus, we can infer that
\begin{equation*}
 {\bf 1}_{\mathcal{C}} \;(a\chi_2'(\phi-\varphi))^{-n} (\theta+{\rm dd}^c u)^n
	\leq {\bf 1}_{\mathcal{C}}(\theta+{\rm dd}^c v)^n \leq {\bf 1}_{\mathcal{C}}(\theta+{\rm dd}^c \varphi)^n,
\end{equation*}
where the last inequality follows from \cite[Lemma 4.5]{DDL5}. The above is equivalent to
\begin{align}
 (\theta+{\rm dd}^c u)^n &={\bf 1}_{ \mathcal{C}} (\theta+{\rm dd}^c u)^n \leq {\bf 1}_{ \mathcal{C}} (a\chi_2'(\phi-\varphi))^{n}(\theta+{\rm dd}^c \f)^n\nonumber\\
 &\leq a^n \chi_1(\phi-\varphi)(\theta+{\rm dd}^c \f)^n, \label{eq: stimabase}
\end{align}
where in the last inequality we used that $(\chi_2')^n=\chi_1$.

We now choose $a$ so that
\[
\beta a^n \int_X \chi_1(\phi-\varphi) (\theta+{\rm dd}^c \varphi)^n=\beta a^n E_{\chi_1}(\varphi, \phi) = m_\phi
\]
 with $\beta > 1$.
Observe that $a\in(0,1)$ since we observed that $E_{\chi_1}(\varphi, \phi)\geq m_\phi$.
Integrating both sides of \eqref{eq: stimabase} over $X$, we obtain
\[
\int_X (\theta+{\rm dd}^c u)^n \leq \frac{m_\phi}{\beta}.
\]
Since $\int_X \theta_u^n=m_\phi$, we arrive at a contradiction.

So we can infer that $G\neq X$, or equivalently that $G^c\neq \varnothing$. This means that there exists~${x_0\in X}$ such that $(\phi-u)(x_0)\leq \tau_2(1)$ where $\tau_2$ is the inverse function of $\gamma_2'$ which we note to be decreasing (since so in $\gamma_2'$). In particular,
\[\sup_X u =\sup_X (u-\phi) \geq -\tau_2(1).\]

Applying the uniform version of Skoda's integrability theorem \cite{Zer01} to $\PSH(X,C\omega)$ for ${C>0}$ such that $\theta\leq C\omega$, we know that there exist uniform constants $c_0, C_0>0$ such that, for all~${h \in \PSH(X,\theta)}$ with $\sup_X h=0$ we have $\int_X {\rm e}^{-c_0 h} \omega^n \leq C_0$. For $h:=u+\tau_2(1)$, we have $\sup_X h \geq 0$, hence
\begin{align*}%\label{key estimate}
 \int_X {\rm e}^{c_0 (a \chi_2(\phi-\varphi) -\tau_2(1))} \omega^n &= \int_X {\rm e}^{c_0 (-\psi+\phi -\tau_2(1))} \omega^n \leq \int_X {\rm e}^{-c_0 (u +\tau_2(1))} \omega^n
\\&= \int_X {\rm e}^{-c_0 (h-\sup_X h) - c_0\sup_X h} \leq C_0,
\end{align*}
where in the first inequality we used $\phi\leq 0$ and $u\leq \psi$.

Let us now give a more explicit expression of $\tau_2(1)$.
By definition, $\tau_2$ is such that
\[s=\tau_2 \left( \frac{1}{a\chi_2'(\gamma_2(s))} \right).\] Now we want to understand $\tau_2(t)$, hence we want to find $s$ such that \smash{$\tfrac{1}{a\chi_2'(\gamma_2(s))}=t $}. This means~${ (at)^{-1}=\chi_2'(\gamma_2(s)) = (\chi_1(\gamma_2(s)))^{1/n} }$ (since $(\chi_2')^n=\chi_1$). Therefore, $\gamma_1((at)^{-n}) =\gamma_2(s)$, where $\gamma_1$ is the inverse function of $\chi_1$. It the follows that $ \tau_2(t)= a\chi_2(\gamma_1((at)^{-n}))$.
In particular, $\tau_2(1)= a\chi_2(\gamma_1(a^{-n}))$.

Also, letting $s:=\gamma_1 (a^{-n})$, we have $a^{-1} = (\chi_1(s))^{\frac{1}{n}}$ so that $a=a(s) = \frac{1}{\chi_2'(s)}$ and
\begin{gather} \tau_2(1) = \frac{\chi_2(s)}{\chi_2'(s)}. \label{Bo} \end{gather}
Note that since $a\in(0,1)$, then $s > 1$ ($\gamma_1$ is increasing and $\gamma(1)=1$).

Next, setting \[K:= \{ x \in X \colon a \chi_2(\phi - \varphi) \leq 2(\phi - \varphi) \},\]
we claim that
\[2^{-1} a \chi_2(\phi - \varphi) \leq a \chi_2(\phi - \varphi) - \tau_2(1) \quad \text{on $X \setminus K$},\]
that is, $a \chi_2(\phi - \varphi) \geq 2 \tau_2(1)$ on $X \setminus K$.

We observe that, since $\chi_2$ is convex and $\chi_2(0) = 0$, the set $E:=\{ t \in \mathbb{R}^+ \colon a\chi_2(t) \leq 2 t \}$ is a~closed convex set containing $0$, i.e., it is an interval of the form $[0, \lambda]$ with $0 \leq \lambda \leq + \infty$. Hence, $x\in X\setminus K$ if and only if $(\phi-\varphi)>\lambda$.
We then need to prove that if $ a\chi_2( \phi - \varphi) > a\chi_2( \lambda)$,
then
$ a \chi_2( \phi - \varphi) \geq {2 \tau_2(1)}$. It is then sufficient to prove that $a\chi_2( \lambda) \geq 2 \tau_2(1)$. The above is equivalent to $ \lambda \geq \gamma_2(2 \tau_2(1)) \geq 0$. By definition of $\lambda$, this means that $\gamma_2(2 \tau_2(1))\in E$, i.e.,
\[\tau_2(1) \leq \gamma_2(2 \tau_2(1)),\qquad \text{or} \qquad a \chi_2(\tau_2(1)) \leq 2 \tau_2(1).\]
By (\ref{Bo}) and $a=1/\chi_2'(s)$, the above inequality gives
\[\chi_2\left(\frac{\chi_2( s)}{\chi_2' (s)} \right) \leq 2 \chi_2(s).\]
Since $\chi_2$ is convex and $\chi_2(0) = 0$, we have
$\chi_2(s) \leq s \chi_2'(s)$.
Then we obtain
\[ \chi_2\left(\frac{\chi_2( s)}{\chi_2' (s)} \right) \leq \chi_2(s) \leq 2 \chi_2(s),\]
and the claim is proved.

It follows that
\begin{align*}%\label{final}
\int_X {\rm e}^{\frac{c_0}{2} a \chi_2(\phi-\varphi)} \omega^n &\leq \int_K {\rm e}^{c_0 (\phi-\varphi)} \omega^n + \int_{X\setminus K}{\rm e}^{c_0(a\chi_2(\phi-\varphi) -\tau_2(1))} \omega^n\\
 &\leq \int_K {\rm e}^{-c_0 \varphi} \omega^n + \int_{X\setminus K}{\rm e}^{c_0(a\chi_2(\phi-\varphi) -\tau_2(1))} \omega^n \\
 &\leq\int_K {\rm e}^{-c_0 (\varphi+1)+c_0} \omega^n + C_0\leq C_0({\rm e}^{c_0}+1),
\end{align*}
wherein the last inequality we applied the uniform Skoda theorem to the function $\varphi + 1$. Recalling that $\beta a^n E_{\chi_1}(\varphi, \phi) = m_\phi$, the result then follows with
$c(\beta) = \frac{c_0}{2} \beta^{(\frac{-1}{n})}m_\phi^{1/n}$, and $C= C_0({\rm e}^{c_0}+1)$.
As $\beta>1$ can be chosen arbitrarily, we observe that since $\phi$ and $\varphi$ are assumed to be equisingular, the function $\exp \big( c(\beta) E_{\chi_1}(\varphi, \phi)^{-\frac{1}{n}} \chi_2(\phi-\varphi)\big)$ is uniformly bounded, so by Lebesgue dominated convergence theorem, we can choose \smash{$c = \frac{c_0}{2} m_\phi^{1/n}$}.
Observe that the constants $c$ and $C$ are independent of $\chi_1$ (and $\chi_2$).
\end{proof}
\begin{thebibliography}{99}
\bibitem[1]{BT87}
Bedford E., Taylor B.A., Fine topology, \v{S}ilov boundary, and {$(dd^c)^n$},
 \href{https://doi.org/10.1016/0022-1236(87)90087-5}{\textit{J.~Funct. Anal.}} \textbf{72} (1987), 225--251.

\bibitem[4]{BBGZ13}
Berman R.J., Boucksom S., Guedj V., Zeriahi A., A variational approach to
 complex {M}onge--{A}mp\`ere equations, \href{https://doi.org/10.1007/s10240-012-0046-6}{\textit{Publ. Math. Inst. Hautes
 \'Etudes Sci.}} \textbf{117} (2013), 179--245, \href{https://arxiv.org/abs/0907.4490}{arXiv:0907.4490}.

\bibitem[7]{BouThesis}
Boucksom S., C\^ones positifs des vari\'et\'es complexes compactes, Ph.D.~Thesis, {U}niversit\'e {J}oseph-{F}ourier-{G}renoble~{I}, 2002.

\bibitem[8]{Bou04}
Boucksom S., Divisorial {Z}ariski decompositions on compact complex manifolds,
 \href{https://doi.org/10.1016/j.ansens.2003.04.002}{\textit{Ann. Sci. \'Ecole Norm. Sup.~(4)}} \textbf{37} (2004), 45--76.

\bibitem[9]{BEGZ10}
Boucksom S., Eyssidieux P., Guedj V., Zeriahi A., Monge--{A}mp\`ere equations
 in big cohomology classes, \href{https://doi.org/10.1007/s11511-010-0054-7}{\textit{Acta Math.}} \textbf{205} (2010), 199--262,
 \href{https://arxiv.org/abs/0812.3674}{arXiv:0812.3674}.

\bibitem[13]{DDL2}
Darvas T., Di~Nezza E., Lu C.H., Monotonicity of nonpluripolar products and
 complex {M}onge--{A}mp\`ere equations with prescribed singularity,
 \href{https://doi.org/10.2140/apde.2018.11.2049}{\textit{Anal. PDE}} \textbf{11} (2018), 2049--2087, \href{https://arxiv.org/abs/1705.05796}{arXiv:1705.05796}.

\bibitem[15]{DDL5}
Darvas T., Di~Nezza E., Lu C.H., The metric geometry of singularity types,
 \href{https://doi.org/10.1515/crelle-2020-0019}{\textit{J.~Reine Angew. Math.}} \textbf{771} (2021), 137--170,
 \href{https://arxiv.org/abs/1909.00839}{arXiv:1909.00839}.

\bibitem[16]{DDL6}
Darvas T., Di~Nezza E., Lu C.H., Relative pluripotential theory on compact
 {K}\"ahler manifolds, \href{https://arxiv.org/abs/2303.11584}{arXiv:2303.11584}.

\bibitem[19]{DGL20}
Di~Nezza E., Guedj V., Lu C.H., Finite entropy vs finite energy,
 \href{https://doi.org/10.4171/cmh/515}{\textit{Comment. Math. Helv.}} \textbf{96} (2021), 389--419,
 \href{https://arxiv.org/abs/2006.07061}{arXiv:2006.07061}.

\bibitem[20]{DNL21}
Di~Nezza E., Lu C.H., Geodesic distance and {M}onge--{A}mp\`ere measures on
 contact sets, \href{https://doi.org/10.1007/s10476-022-0159-1}{\textit{Anal. Math.}} \textbf{48} (2022), 451--488,
 \href{https://arxiv.org/abs/2112.09627}{arXiv:2112.09627}.

\bibitem[21]{DNT23}
Di~Nezza E., Trapani S., The regularity of envelopes, \textit{Ann. Sci. \'Ec.
 Norm. Sup\'er.}, {t}o appear, \href{https://arxiv.org/abs/2110.14314}{arXiv:2110.14314}.

\bibitem[24]{GZ07}
Guedj V., Zeriahi A., The weighted {M}onge--{A}mp\`ere energy of
 quasiplurisubharmonic functions, \href{https://doi.org/10.1016/j.jfa.2007.04.018}{\textit{J.~Funct. Anal.}} \textbf{250}
 (2007), 442--482, \href{https://arxiv.org/abs/math.CV/0612630}{arXiv:math.CV/0612630}.

\bibitem[26]{Gup23}
Gupta P., A complete metric topology on relative low energy spaces,
 \href{https://doi.org/10.1007/s00209-023-03218-5}{\textit{Math.~Z.}} \textbf{303} (2023), 56, 27~pages, \href{https://arxiv.org/abs/2206.03999}{arXiv:2206.03999}.

\bibitem[31]{Zer01}
Zeriahi A., Volume and capacity of sublevel sets of a {L}elong class of
 plurisubharmonic functions, \href{https://doi.org/10.1512/iumj.2001.50.2062}{\textit{Indiana Univ. Math.~J.}} \textbf{50}
 (2001), 671--703.

\end{thebibliography}
\end{document}
